\documentclass[11pt]{article}
\usepackage[margin=1in]{geometry}
\usepackage[T1]{fontenc}
\usepackage{hyperref}
\usepackage{parskip}
\usepackage{setspace}
\usepackage{titlesec}

\titleformat{\section}{\large\bfseries}{}{0em}{}[\titlerule]
\titlespacing{\section}{0pt}{1.5em}{0.75em}

\hypersetup{colorlinks=true, urlcolor=blue, linkcolor=black}

\title{\textbf{Semantic Odyssey: Navigating Word Similarity Graphs}}
\author{
    Alexander Davydenko \and
    Sophia LAST NAME \and
    Abhirve LAST NAME \and
}
\date{CSC111 Project 2 --- Phase 1 Proposal \\ March 2026}

\begin{document}

\maketitle

% ─────────────────────────────────────────────────────────────────────────────
\section{Problem Description and Project Question}

Human languages are vast, and the meanings of words exist on a continuous
spectrum of similarity. A word like ``fire'' shares something with ``flame,''
which in turn relates to ``heat,'' which connects to ``sun,'' and one can
plausibly trace a path all the way to ``ocean.'' The study of how words relate
semantically is a central topic in computational linguistics, and modern natural
language processing has produced powerful mathematical tools for quantifying
these relationships.

One such tool is GloVe (Global Vectors for Word Representation), a set of
pre-trained word embeddings produced by Pennington et al.\ at Stanford
\cite{glove}. In GloVe, each word is represented as a vector in
high-dimensional space; words that tend to appear in similar linguistic contexts
are placed near each other in that space. The \emph{cosine similarity} between
two word vectors gives a principled, numerical measure of their semantic
closeness, ranging from $-1$ (opposites) to $+1$ (identical context). This
structure naturally gives rise to a \emph{graph}: the vertices are words, and
edges connect words whose cosine similarity exceeds a threshold, weighted by
that similarity score.

The games Semantle and Wordle demonstrated widespread public interest in
word-meaning puzzles. Inspired by these, we designed \emph{Semantic Odyssey}:
an interactive word-navigation puzzle in which the player starts at one word and
must navigate a semantic similarity graph---step by step, choosing from a set of
nearby words---to reach a hidden target word in as few moves as possible. The
player receives feedback on how ``close'' or ``distant'' each available word is,
and the game computes a provably achievable optimal path using breadth-first
search on the same graph the player navigates.

\textbf{Our project question is: How can the semantic similarity graph induced
by pre-trained GloVe word embeddings be exploited to build an engaging,
challenge-calibrated word-navigation game, and how can BFS on this implicit
graph provide optimal-path benchmarks that are provably achievable by a human
player navigating the same graph?}

This question is interesting because the graph is not precomputed but
constructed lazily on demand from a 400,000-word embedding model---making the
graph exploration both computationally rich and grounded in real linguistic data.


% ─────────────────────────────────────────────────────────────────────────────
\section{Computational Plan}

\subsection*{Data}

Our primary dataset is the GloVe \texttt{wiki-gigaword-50} corpus: 400,000
English word vectors, each of dimension 50, trained on approximately 6 billion
tokens drawn from Wikipedia 2014 and the English Gigaword 5 corpus \cite{glove}.
We access the dataset through the \texttt{gensim} library's model downloader
(\texttt{gensim.downloader.load}), which fetches and caches the vectors
automatically. A sample from the data: the five nearest words to ``fire'' by
cosine similarity are ``fires'' (0.82), ``wildfire'' (0.80), ``blaze'' (0.79),
``flames'' (0.77), and ``burning'' (0.75), illustrating the rich semantic
neighbourhood structure the vectors encode.

\subsection*{Graph Representation}

We model the semantic word network as a directed, weighted, \emph{implicit}
graph $G = (V, E)$ where $V$ is the GloVe vocabulary.  Rather than
materialising all $|V|^2 \approx 1.6 \times 10^{11}$ possible edges, we
construct the graph lazily: for any vertex $w \in V$, its out-neighbourhood
$N(w)$ is computed on demand by querying
\texttt{model.most\_similar(w, topn=50)}, returning the 50 words with the
highest cosine similarity to $w$.

From these 50 candidates we apply a \emph{fixed-bucket selection} strategy to
ensure players always see a diverse spread of options:
\begin{itemize}
    \item 3 \emph{close} words with cosine similarity $\geq 0.65$,
    \item 4 \emph{related} words with similarity in $[0.45,\,0.65)$,
    \item 3 \emph{distant} words with similarity $< 0.45$,
\end{itemize}
yielding a consistent out-degree of 10 for each vertex.  If a bucket contains
fewer than its quota the shortfall is filled from the remaining candidates in
descending similarity order.  All computed neighbourhoods are memoised in a
dictionary keyed by word so each GloVe query is performed at most once per
session.

\subsection*{Computations}

\textbf{Breadth-first search (BFS).}
Our core computation is BFS on the implicit graph.  Given a start word $s$ and a
target word $t$, BFS expands nodes level by level---each node's successors being
its pre-computed out-neighbourhood---and returns the shortest path from $s$ to
$t$ in terms of number of steps.  Because the neighbour list is deterministic
and fixed, the BFS path is provably achievable by a human player following the
same edges; the shown ``optimal path'' is never a fiction.

\textbf{Puzzle generation.}
Puzzle generation applies BFS with a fixed daily seed to identify suitable
start-target pairs.  We prefer pairs whose shortest BFS path is between 4 and 8
steps, balancing challenge and solvability.  We run a fast, target-unaware BFS
during the loading stage to validate candidates quickly (sharing a cache across
all calls), then compute the full target-aware optimal path in a background
thread concurrently with gameplay.

\textbf{Target guarantee.}
At each player step we compute the cosine similarity between the current word
and the target.  If this similarity is $\geq 0.60$---a threshold that captures
morphological variants (e.g.\ ``moons'' $\to$ ``moon,'' ``roses'' $\to$
``rose'') while excluding merely related pairs---the target is inserted into the
10-word neighbour list regardless of bucket placement, ensuring the player can
always take the final step once they are genuinely close in semantic space.

\textbf{Path comparison.}
On completing the puzzle the player's step count is compared against the BFS
optimal, and the full paths are rendered side by side so the player can study
where they diverged.

\subsection*{Visualisation and Interaction}

The program is a Pygame application with four screens managed by a state
machine: (1) a \emph{loading screen} with an animated progress bar that
reflects staged progress through vector loading and puzzle generation;
(2) a \emph{start screen} showing the target word and three selectable starting
words; (3) the \emph{game screen} showing the current word, the player's path so
far, and a $2 \times 5$ grid of clickable word tiles colour-coded by similarity
label (\emph{close}, \emph{related}, \emph{distant}); and (4) a \emph{result
screen} displaying the player's path against the BFS-optimal path with step
counts.  An undo button allows players to retrace steps.  After winning, the
player may continue to a new randomly chosen target from their current word.

\subsection*{New Library: \texttt{gensim}}

We use the \texttt{gensim} library (version 4.x) \cite{gensim} specifically for
its pre-trained word vector interface, which is not covered in CSC110/111.  Key
components we rely on are:
\begin{itemize}
    \item \texttt{gensim.downloader.load('glove-wiki-gigaword-50')} --- fetches
          and caches the 400,000-word embedding model;
    \item \texttt{KeyedVectors.most\_similar(word, topn=k)} --- retrieves the
          $k$ nearest neighbours of a word by cosine similarity in
          $O(|V| \cdot d)$ time ($d = 50$), which is the engine of our implicit
          graph construction;
    \item \texttt{KeyedVectors.similarity(w1, w2)} --- computes pairwise cosine
          similarity in $O(d)$ time, used for the target guarantee check at
          every player step.
\end{itemize}
Without \texttt{gensim}, computing cosine similarities over a 400,000-word
vocabulary from scratch would require either prohibitive preprocessing or a
custom embedding, making it infeasible within the project scope.  The library is
appropriate because it provides exactly the semantic graph structure our project
requires while leaving all graph traversal, puzzle logic, and game rendering to
our own code.


% ─────────────────────────────────────────────────────────────────────────────
\section*{References}

\begin{thebibliography}{9}

\bibitem{glove}
    J. Pennington, R. Socher, and C. D. Manning,
    ``GloVe: Global Vectors for Word Representation,''
    in \textit{Proceedings of the 2014 Conference on Empirical Methods in
    Natural Language Processing (EMNLP)}, Doha, Qatar, 2014, pp.\ 1532--1543.
    [Online]. Available: \url{https://nlp.stanford.edu/pubs/glove.pdf}

\bibitem{glovedataset}
    Stanford NLP Group,
    ``GloVe: Global Vectors for Word Representation --- Pre-trained vectors,''
    2014.
    [Online]. Available: \url{https://nlp.stanford.edu/projects/glove/}

\bibitem{gensim}
    R. \v{R}eh\r{u}\v{r}ek and P. Sojka,
    ``Software Framework for Topic Modelling with Large Corpora,''
    in \textit{Proceedings of the LREC 2010 Workshop on New Challenges for NLP
    Frameworks}, Valletta, Malta, 2010, pp.\ 45--50.
    [Online]. Available: \url{https://radimrehurek.com/gensim/}

\bibitem{gensim-docs}
    gensim Development Team,
    \textit{gensim 4.x Documentation: models.keyedvectors},
    2023.
    [Online]. Available:
    \url{https://radimrehurek.com/gensim/models/keyedvectors.html}

\bibitem{pygame}
    Pygame Community,
    \textit{Pygame Documentation},
    2023.
    [Online]. Available: \url{https://www.pygame.org/docs/}

\bibitem{semantle}
    D. Turner-Evans,
    ``Semantle,''
    2022.
    [Online]. Available: \url{https://semantle.com/}

\end{thebibliography}

\end{document}
